\documentclass[11pt]{article}

% --------------------------------------------------
% Packages
% --------------------------------------------------
\usepackage[T1]{fontenc}
\usepackage{lmodern}
\usepackage{geometry}
\usepackage{microtype}
\usepackage{setspace}
\usepackage{amsmath,amssymb,amsthm}
\usepackage{csquotes}
\usepackage{hyperref}

\geometry{margin=1in}
\setstretch{1.15}

% --------------------------------------------------
% Theorem Environments
% --------------------------------------------------
\newtheorem{definition}{Definition}
\newtheorem{proposition}{Proposition}
\newtheorem{remark}{Remark}

% --------------------------------------------------
% Macros
% --------------------------------------------------
\newcommand{\Fut}{\mathcal{F}}
\newcommand{\State}{\Sigma}
\newcommand{\Event}{\mathcal{E}}
\newcommand{\Gen}{\mathcal{G}}

% --------------------------------------------------
% Title
% --------------------------------------------------
\title{Event-Nodes and Texture-Nodes:\\
Extending Causal Mechanistic Interpretability with Event-Historical Structure}
\author{Flyxion}
\date{December 2025}

\begin{document}
\maketitle

% ==================================================
\begin{abstract}
% ==================================================

Recent work in mechanistic interpretability has advanced a rigorously causal
framework for understanding neural networks by treating internal components as
variables within an interventional causal model. While this paradigm succeeds
in distinguishing correlation from mechanism and enables algorithmic
abstraction via counterfactual testing, it remains implicitly state-centric and
only weakly sensitive to irreversibility and history.

This essay introduces a minimal extension to the causal--mechanistic framework
by distinguishing two types of causal nodes: \emph{texture-nodes}, which carry
instantaneous values and are exhaustively characterized by value substitution,
and \emph{event-nodes}, which record irreversible structural commitments that
constrain the space of future admissible states. Using a procedural generation
analogy, texture-nodes correspond to parameter values of a generator, while
event-nodes correspond to the selection of the generator itself.

We show that this distinction preserves the full rigor of causal abstraction
while explaining substitution failure, semantic incompatibility, algorithmic
mixtures, and structural irreversibility without invoking metaphysical
assumptions. Mechanistic understanding is reframed as the reconstruction of both
value-level mediation and event-level constraint structure under intervention.

\end{abstract}

\tableofcontents

% ==================================================
\section{Introduction}
% ==================================================

Mechanistic interpretability seeks to answer a precise question: what internal
structures of a neural network are responsible for its behavior, and in what
sense? A purely observational approach—examining activations, attention maps,
or weight magnitudes—fails to answer this question, as correlation alone cannot
establish causal responsibility.

Recent work has therefore adopted a causal paradigm, treating neural networks
as systems amenable to intervention, counterfactual reasoning, and algorithmic
abstraction. Under this view, understanding a model mechanistically requires
identifying internal components whose manipulation predictably alters behavior
in ways consistent with a hypothesized computation.

However, this paradigm remains implicitly committed to a state-based ontology:
interventions act on instantaneous variable values, and history matters only
insofar as it affects present behavior. This essay argues that such a framework
is incomplete. Certain causal phenomena—irreversibility, commitment,
substitution failure, and algorithmic regime shifts—cannot be fully explained
without treating events themselves as first-class causal objects.

We propose a minimal extension: a typed causal graph distinguishing
\emph{texture-nodes} from \emph{event-nodes}.

% ==================================================
\section{Causal Mechanistic Interpretability}
% ==================================================

In the narrow technical sense, mechanistic interpretability is defined by causal
explanatory adequacy. A component is mechanistically relevant if intervening on
it changes model behavior in a systematic and predictable way.

Interventions are typically implemented as value substitutions: fixing an
internal activation to a value it would not naturally assume, or replacing it
with the value observed under a different input. Counterfactual success is
measured by alignment between low-level interventions and high-level algorithmic
predictions.

This methodology supports \emph{causal abstraction}: if interventions on
low-level variables produce the same counterfactual outcomes as interventions
on a high-level algorithm, the algorithm is considered an abstraction of the
network.

Despite its power, this framework leaves open questions concerning history,
irreversibility, and the conditions under which substitution itself is well
defined.

% ==================================================
\section{Texture-Nodes: Value-Carrying Variables}
% ==================================================

\begin{definition}[Texture-Node]
A \emph{texture-node} is a causal variable whose role is exhaustively determined
by the value it instantiates at a given time.
\end{definition}

Formally, a texture-node corresponds to a variable \( x_t \in \State \), where
causal influence is evaluated by substituting \( x_t \mapsto x_t' \) and
observing downstream effects.

Texture-nodes include neurons, attention heads, MLP activations, and linear
directions used in activation steering. Their causal relevance is established
through value-level interventions and mediation analysis.

Texture-nodes admit reversible manipulation: if the original value can be
restored, the system returns to its prior functional state.

% ==================================================
\section{Event-Nodes: Constraint-Carrying Commitments}
% ==================================================

\begin{definition}[Event-Node]
An \emph{event-node} is a causal variable whose role is not reducible to an
instantaneous value, but consists in imposing a structural constraint on the
space of admissible future states.
\end{definition}

An event-node induces a transformation
\[
\Fut_{t+1} \subseteq \Fut_t,
\]
where \( \Fut_t \) denotes the set of reachable futures from time \( t \).

Unlike texture-nodes, event-nodes are not exhaustively characterized by value
substitution. Their causal influence is evaluated by changes in what the system
can or cannot subsequently do.

Event-nodes are generally irreversible: no sequence of later texture-node
interventions can restore the original future space once a structural
commitment has been made.

% ==================================================
\section{Generator Functions and Procedural Structure}
% ==================================================

The distinction between event-nodes and texture-nodes can be clarified through a
procedural generation analogy.

Let \( \Gen \) be a family of generator functions, each parameterized by a vector
\( \theta \in \mathbb{R}^n \). A generator \( G \in \Gen \) defines a family of
outputs
\[
y = G(\theta).
\]

Texture-nodes correspond to the parameters \( \theta \) or to ratios
\( \theta_i / \theta_j \), which function as vector-like directions in parameter
space. Adjusting these values modifies surface appearance without changing the
generative regime.

Event-nodes correspond to the selection of the generator \( G \) itself. Once a
generator is chosen, all admissible parameterizations are conditioned on that
choice. No amount of parameter tuning reproduces the output family of a
different generator.

This distinction explains why some interventions modulate behavior smoothly,
while others produce qualitative regime shifts.

% ==================================================
\section{Irreversibility and History}
% ==================================================

Standard causal analysis treats irreversibility as a consequence of information
loss. Event-nodes introduce a second, more fundamental notion.

\begin{definition}[Structural Irreversibility]
An intervention is \emph{structurally irreversible} if it reduces the space of
admissible futures in a way that cannot be undone by subsequent value-level
interventions.
\end{definition}

Structural irreversibility does not depend on signal destruction. It arises from
commitment to a generative structure rather than from noise or ablation.

This explains why two systems with identical current states may behave
differently under future interventions: they may share texture-state while
differing in event-history.

% ==================================================
\section{Causal Abstraction with Typed Nodes}
% ==================================================

A causal abstraction must now satisfy two criteria:

\begin{enumerate}
\item \textbf{Texture Correspondence}: value-level interventions align between
the abstraction and the underlying system.
\item \textbf{Event Correspondence}: structural commitments and irreversibility
relations are preserved.
\end{enumerate}

Multiple abstractions may remain valid at different levels, but abstractions that
erase event-level structure mischaracterize the system’s causal organization,
even if short-horizon input--output behavior is preserved.

Algorithmic mixtures and regime switching are naturally interpreted as
event-node transitions rather than as noise or failure.

% ==================================================
\section{Failure, Substitution, and Semantic Incompatibility}
% ==================================================

Interchange interventions assume semantic compatibility between substituted
values and their new contexts. Event-nodes explain why this assumption sometimes
fails.

Substitution failure does not indicate a mysterious semantic breakdown; it
indicates that the substituted value was generated under a different event
history and is incompatible with the current future constraints.

Novel behavior under intervention signals an incorrect or incomplete causal
graph, often due to unmodeled event-nodes.

% ==================================================
\section{What Counts as Understanding}
% ==================================================

Mechanistic understanding is not merely the ability to control or predict
behavior. It consists in reconstructing:

\begin{itemize}
\item how values mediate effects (texture-level),
\item which commitments constrain futures (event-level),
\item and how these interact under intervention.
\end{itemize}

Understanding is necessarily partial and revisable. No single abstraction
exhausts the system’s causal structure, but some abstractions fail by erasing
irreversibility and history.

% ==================================================
\section{Conclusion}
% ==================================================

By distinguishing texture-nodes from event-nodes, we obtain a minimal extension
to causal mechanistic interpretability that preserves its empirical rigor while
accounting for irreversibility, history sensitivity, and structural commitment.

This extension requires no new metaphysics and no abandonment of causal
abstraction. It merely clarifies what interventions act upon and what kinds of
causal structure must be preserved for an explanation to remain valid.

Mechanistic interpretability, properly extended, becomes not only the study of
what internal variables do, but of what systems have irreversibly become.

\end{document}
